\section{The family bound at primes dividing the small conductor and at three}

% BEGIN G82F1
\paragraph{G82F1. The split declaration and the cyclotomic level.}
Use the member, setup and representation of F78D2. Fix a twisted
localization $i$, a good point $\pi$, and $p=\ell_i$ with
$p\in D.\mathrm{host.small}$ or $p=3$. Write $Q_S$ for the
host's small conductor, $N=Q_S qd$, and
$s=\max(1,v_p(N))$. The small set is the set of prime divisors
of $Q_S$, which is prime to $qd$. It includes $2$, and in fact
$4\mid Q_S$ by the host tower divisor. The host conditions give
$q,d\ne3$ and $q\ge11$. Neither type prime can occur: each is
$1\bmod N$ and none of the prime divisors of $Q_S$ is $1\bmod N$,
while $3<N$. Thus $v_p(N)=v_p(Q_S)$ for every prime now tested.
The native exponent $s_0(i)$ at an odd prime is
$\max(1,v_p(Q_S))$, and at $2$ is $v_2(Q_S)\ge2$. Hence it is $s$.

At an odd $p\mid Q_S$ the host condition on divisors of $Q_S$
gives $(-d/p)=1$. At $p=3$ outside $Q_S$ it follows from the
host small-prime splitting condition because the setup's split
bound is at least $60$. Hensel's lemma gives a square root of
$-d$ in $\mathbb Q_p$. At $2$ the host gives $d\equiv7\bmod8$:
for $f(X)=X^2+d$ the integer $1$ satisfies
$v_2(f(1))\ge3>2v_2(f'(1))$, so a root exists in $\mathbb Z_2$
as well. In every case both distinct roots lie in $\mathbb Q_p$.
It follows that all embeddings at one place give the same value
on $\sqrt{-d}$. Its conjugate place gives the opposite value:
the prime of $\mathbb Q(\sqrt{-d})$ below it is the conjugate
one of the two split primes. This also distinguishes the two
places at $2$, independently of reduction of the two roots.
Use the orientation from S80J1 and declare its upward places
columns and their conjugates rows, as in G81S1. These sets
partition the places above $p$. The local degrees $d_u$ and $d_w$
on a conjugate pair agree, including when the host completion is
ramified. The sum of all those degrees is $[E:\mathbb Q]=2f=4n$.

The bridge completion $K=E_w$ lies in
$\mathbb Q_p(\mu_{\operatorname{lcm}(p,N)})$ by the host
cyclotomic embedding: its embedded dense subfield is generated
inside the conductor cyclotomic field, and the local cyclotomic
field is finite and closed. The good coefficient extension contains
these roots. This proves the local $\mathrm{Cyc}$ hypotheses,
including the dyadic level. At each such place inertia contains
an element with cyclotomic value $1+p^s$. Indeed the cyclotomic
$p$-power tower is totally ramified, so local inertia over
$\mathbb Q_p$ surjects onto its Galois group, by finite-level
surjectivity and compactness. An inertia lift of $1+p^s$ fixes
all prime-to-$p$ roots of unity by unramifiedness, as well as
$\mu_{p^s}$ by its cyclotomic value. Thus it fixes $K$ and gives
an inertia element there, transported through the bridge.
For $j>0$ one has
\[
      v_p((1+p^s)^j-1)=s+v_p(j).
\]
The first binomial term has strictly smallest valuation for an
exponent prime to $p$. Raising successively to the $p$th power
adds one to the valuation, using $s\ge2$ when $p=2$. Inverting
does not change the valuation. This bounds the normalized length
of $(L/O)(-j)^{D_u}$ by $s+v_p(j)$, since a fixed element is
killed by this scalar difference. Thus the column maximum
$h_0$ is at most $s+\lfloor\log_p t\rfloor$.
% END G82F1

% BEGIN G82F2
\paragraph{G82F2. The augmented determinant with the small-prime exponents.}
Write $h_B=b+S_c+\eta+t+2$ and use the three heights
$h_1=\lfloor h_B\rfloor$, $h_2=\lfloor2b\rfloor+1$,
$h_3=\lfloor2b\rfloor$. Put
\[
 \begin{gathered}
 a=e_L(p,N,h_1),\qquad c_0=c_{0L}(p,N,h_2,4n n_V),\\
 c_B=c_{BL}(p,N,h_3,4n n_V,t).
 \end{gathered}
\]
These are the constants in the original small-prime split bound.
We verify the changes of hypotheses needed for G82C1--G82C5,
whose initial scope was the unramified middle range.

The good point implies $\pi\in Z_i$ and nonempty good support.
D79G1--D79G3 give the Hecke characteristic polynomials, the
full integral rank block and the two projectors. G81S3's
normalization applies at $\pi$. Use $g_0$ given by G82F1, so its
$t$ values there are $\chi_1(g_0)(1+p^s)^{-j}$, $0\le j<t$.
Away from $\pi$ the corner perturbation of G81S3 applies here
too. Every localized full representation is crystalline by S80J3
since $p$ differs from the type primes. Its labelled jumps after
the common twist are simple by S80J1. Strictness therefore excludes
repeated isomorphic constituents, exactly the hypothesis used to
make the corner discriminants nonzero. The supported perturbations
preserve the matrix at $\pi$. After a common finite enlargement
we get the free ordered-root algebra and $E^2=\delta E$ with
nonzero $\delta$ at every lift. At the distinguished lift,
\[
            v=v_p(\delta(\pi'))=\sum_{1\le j<t}(s+v_p(j))
                                 =D.\mathrm{nd.vdel}(i).
\]
The factor $\chi_1(g_0)$ is a unit. D79G7 supplies the global
off-support condition. Its auxiliary polynomial and quadratic
identities hold for these same primes, including nonvanishing
of $\zeta_P-1$ at $P\in Q_i$. The unit test at the designated
order-$15$ auxiliary character in $S_i$ follows from
$\Phi_{15}(1)=1$, including when $p\mid15$.

At each tested coefficient extension S80J1--S80J3 place the full
twisted family and its dual twist in $[-h_1,h_1]$, since the
original jumps have absolute value at most $h_B-2$ and the dual
changes $z$ to $-1-z$. Strict crystalline subquotients inherit
that interval. The two onto left-ideal evaluations of F82L1,
and the dual maps on the graph lattices of I81G3, are therefore
covered by L82K5 over the fields of G82F1. This supplies every
lift in F82L1--F82L3 at $a$. The column comparison of F82L4
uses the nonnegative character string $0,\dots,t-1$ and the
crystalline dimension formula. Both are valid on these completions.
Moreover the rank cloud is positive on an upward column and the
dual rank cloud is positive on its downward row. Each rational
finite rank is at most $d_u m$, and each relevant rational
invariant space vanishes, by the jump argument of G81S1.

On a column for $W$ or its conjugate row for $W^\vee(1)$,
zero is absent and the jumps lie in $[-h_2,h_2]$ by S80J2.
The union of the jumps has at most $4n n_V$ elements, since
there are at most $4n$ local embeddings and at most $n_V$
jumps per embedding. T82K6 kills the divisible invariants by
$p^{c_0}$, including over the extended coefficients.
Consequently $h_W\le r c_0$ and the dual row cokernel has
normalized length at most $d_um a+r c_0$: as in G81S2 its
finite integral target has free rank at most $d_um$ and torsion
of normalized length at most $r c_0$. On the ordinary ambient
family all jumps have absolute value at most $h_3$ by S80J2,
and $t\le h_3$ from the size bounds. At each place the common
jump lists, independent of the point, have union of size at most
$4n n_V$. The character cut of I81G5 is thus covered by
T82K6's torsion bound $c_B$. The cyclic obstruction proof
of M81B3 uses G82F1's element $1+p^s$ in place of $1+p$:
its $(j+1)$st power differs from one, so the dual invariant
group is cyclic and rational $H^2$ vanishes. The local Euler and
finite-part dimensions in that proof give the same bound
$(d_w+1)c_B$ at every row, including in residue characteristic two.

We can now use the presentation and determinant proof of
G82C1--G82C4, with these local cocycles and constants. The same
rank-nullity argument gives a nonzero evaluation on each point
over the complement, and the determinant polynomials descend
without a residue-cardinality condition. The rank classes are
finite on all columns, the nontrivial character classes vanish
locally there, and the first-character classes are unramified
there. At the other places use G79E6, G79E8 and I81G3--I81G6
as in G82C5. In particular the native killing condition bounds
the auxiliary loss $t\sum_{P\in Q_i}z_P/k_O$ by $p^{-2}$.
The original $\mathrm{TwistedPlaces.ShaFinOdd}$ in the family
setup quantifies over all primes, including $2$, for every
$0<j<t$ (Y79T4). Thus the Selmer length given by killing and
all the required cyclotomic lengths are finite. I81G2 and the
evaluation-ideal argument in G81S3 descend all normalized lengths
to the original DVR.

Writing $c$ for the number of columns, this gives
$\mathrm{Bd}(X.\rho,i,\pi,\mathcal E)$ with
\[
 \begin{aligned}
 \mathcal E={}&
   \sum_{u\in\mathrm{col}}\{(2d_um+1)a+d_um v+2v+h_W(i,u,\pi)+r c_0\}\\
   &+t c_B\sum_{w\in\mathrm{row}}(d_w+2)
                         +4 n_V(c+1)v+tc h_0 .
 \end{aligned}
\]
Indeed $\mathrm{Bd}$ already includes $s_\pi$, the cyclotomic
sum over $0<j<t$, $v_p(h_E)$ and the $p^{-2}$ allowance.
This construction uses a declaration chosen for this index on
the given representation. $\mathrm{StarF}$ does not prescribe
the hub's declaration at these primes.
% END G82F2

% BEGIN G82F3
\paragraph{G82F3. The split error is bounded by the two small-prime costs.}
Put $y=h_B$, $\Lambda=s+\log_p(2y)+1$. For a sufficiently
late member $b\ge1$, $S_c+\eta+t\le b$ by S80J2's parameter
bounds, so $b\le y\le6b$ and all $h_i\le4y$. Also $N\le b$,
since $N^{36}t^{150}\le b^{1/2}$ by the conductor condition
and $N,t\ge1$. Let
$\lambda_N(h)=s+\lfloor\log_p(2h)\rfloor+
\lfloor\log_p N\rfloor+1$.
For $h\le4y$ both integer logarithms are at most
$\log_p(2y)+2\le\Lambda$. Hence $\lambda_N(h)\le3\Lambda$.
Also $\Lambda\le6ys$: use $p\ge2$, $s,y\ge1$ and
$\log(2y)\le2y-1$, $\log2>1/2$. From the definitions
\[
 e_L=32(2hs+1)+\lambda_N(h),\quad
 c_{0L}=20(n'+1)\lambda_N(h),\quad
 c_{BL}=40(t+1)(n'+1)\lambda_N(h)
\]
with $n'=4n n_V=2f n_V$, we obtain, with $D_0=1000$,
\[
  a\le D_0ys,\quad c_0\le D_0 f n_V\Lambda,\quad
  c_B\le D_0 t f n_V\Lambda,\quad h_0\le D_0\Lambda.
\]
For example $n'+1\le3f n_V$ and $t+1\le2t$, giving the
coefficient $720$ in the third bound. For the last bound use
G82F1 and $t\le y$. Further
$v\le t(s+1)\le t\Lambda$, since
$\sum_{1\le j<t}v_p(j)=v_p((t-1)!)\le t$ by the sum of
the floors $(t-1)/p^k$.

Both the number of places and the sum of their local degrees
on either side of the declaration are at most $2f$ by G82F1.
The terms of $\mathcal E$ involving $a$ therefore sum to at
most $6D_0 f m y s$, using $m\ge1$. For the remaining terms
use $m\le f$, $r\le t$, $h_W\le r c_0$,
and $f,t,n_V\ge1$. In units of $D_0 f^2 t^2 n_V\Lambda$,
the $\sum_u d_um v$ terms cost at most $2$, the
$\sum_u 2v$ terms at most $4$, and the two rank-invariant
terms together at most $4$. Next $\sum_w(d_w+2)\le6f$,
so the row characters cost at most $6$. Finally $c+1\le3f$
gives cost at most $12$ for the term $4 n_V(c+1)v$, and
$tc h_0$ costs at most $2$. Summing gives the stated pointwise
cost
\[
      \mathcal E\le 6000 f m y s+30000 f^2t^2 n_V\Lambda .
\]
% END G82F3

% BEGIN G82F4
\paragraph{G82F4. One family constant and a table for the entire small-prime field.}
Let $C_1\ge1$ be the family constant of A78P3 and take $x$
beyond its cutoff with $\log x\ge1$. We have $n,t\le x^{C_1}$
and $f\le k\le r\le t$. The host bound $Q_S\le n^2$ gives
$2^{v_p(Q_S)}\le p^{v_p(Q_S)}\le n^2$. Consequently
$s=\max(1,v_p(Q_S))\le3C_1\log x$, since $\log2>2/3$.
Putting $U=L.\mathrm{sc}.L=\log b\ge1$ and using $y\le6b$
gives $\log_p(2y)\le2U+6$. Therefore
$\Lambda\le(3C_1+9)xU$, since $x\ge1$, $\log x\le x$
and $U\ge1$. Also $f^2t^2 n_V\le2x^{5C_1}$ because
$n_V=r+t\le2t$. For any common constant
\[
 C\ge 18\cdot6000 C_1+2\cdot30000(3C_1+9)+5C_1+1
\]
the first term of G82F3 is at most $C f m b\log x$;
the second is at most $C x^C U^2$ since $C\ge5C_1+1$.
Their sum is the native $\mathrm{midF}(L,C)$ allowance.
Choose $C$ also above the previous family lower bounds and take
the common cutoff for this choice before selecting the member
and hub, as in Y79T4 and J81H2.

At that member define a native auxiliary table $t_f$ by
$t_f.\mathrm{len}(i,w)=0$ and
$t_f.\mathrm{BB}(i)=\mathrm{midF}(L,C)$ for every $i,w$.
Give its $s$, $\mathrm{gcost}$ and $\mathrm{gs}$ entries value
zero, its band the empty predicate, and its exceptional set the
empty set. On every tested row the sum of the local lengths is
zero. Adding the non-place term gives equality in the table
condition $\mathrm{SmallPrimeLengthBound}(L,t_f,C)$. For every twisted $i$ with $p$ small or
$p=3$ and every $\pi\in Z_{\rm good}(i)$, G82F2--G82F3
give the required $\mathrm{Bd}$ with this table sum. This
constructs the existential table in $\mathrm{StarF}(L,X.\rho,C)$
and proves the original field $\mathrm{Concl.two}$ on the same
family hub for every late member with its original finiteness
and scale guards.
% END G82F4
